\documentclass[a4paper,12pt]{article}
\usepackage[utf8]{inputenc}
\usepackage[T1]{fontenc}
\usepackage[french]{babel}
\usepackage{geometry}
\usepackage{graphicx}
\usepackage{hyperref}
\usepackage{enumitem}
\usepackage{xcolor}

\geometry{margin=2.5cm}

\title{\textbf{Rapport d'avancement - Projet NECFORM}}
\author{[Votre nom]}
\date{[Date du jour] \\ Encadrant : [Nom de l'encadrant]}

\begin{document}

\maketitle

\section*{Présentation du projet}

NECFORM est une application web de gestion des formations destinée à faciliter l'organisation et le suivi des formations professionnelles.

L'application permet de :
\begin{itemize}
    \item Gérer les utilisateurs (apprenants, formateurs, administrateurs)
    \item Gérer les entreprises clientes
    \item Gérer le catalogue de formations
    \item Planifier et gérer les sessions de formation
    \item Gérer les inscriptions aux sessions
    \item Traiter les demandes de nouvelles formations
    \item Gérer les documents associés
    \item Visualiser un tableau de bord de synthèse
\end{itemize}

\section{Fonctionnalités développées}

\subsection{Gestion des utilisateurs}
\begin{itemize}
    \item Création, modification et suppression d'utilisateurs
    \item Gestion des rôles et des profils
    \item Association aux entreprises
\end{itemize}

\subsection{Gestion des entreprises}
\begin{itemize}
    \item Création et gestion des entreprises clientes
    \item Association des utilisateurs aux entreprises
\end{itemize}

\subsection{Gestion des formations}
\begin{itemize}
    \item Catalogue de formations avec description
    \item Création, modification et suppression de formations
    \item Gestion des types de formations
\end{itemize}

\subsection{Gestion des sessions}
\begin{itemize}
    \item Planification des sessions de formation
    \item Gestion des dates et lieux
    \item Suivi des statuts (planifiée, en cours, terminée, annulée)
\end{itemize}

\subsection{Gestion des inscriptions}
\begin{itemize}
    \item Inscription des utilisateurs aux sessions
    \item Suivi des statuts d'inscription (en attente, confirmée, annulée)
\end{itemize}

\subsection{Demandes de formation}
\begin{itemize}
    \item Système de demande de nouvelles formations
    \item Workflow de validation des demandes
    \item Suivi des statuts (en attente, approuvée, refusée)
\end{itemize}

\subsection{Gestion des documents}
\begin{itemize}
    \item Upload et gestion des documents
    \item Classification par type (contrat, facture, attestation, etc.)
\end{itemize}

\subsection{Dashboard}
\begin{itemize}
    \item Vue synthétique de l'activité
    \item Statistiques globales
\end{itemize}

\section{Modèle de données et relations entre entités}

\subsection{Vue d'ensemble des relations}

Le système de gestion des formations repose sur les relations suivantes entre les entités principales :

\begin{verbatim}
Entreprise ──(1:N)──> DemandeFormation ──(N:1)──> Formation
                                              │
                                              │ (1:N)
                                              ↓
                                     SessionFormation ──(1:N)──> Inscription
                                              │                      │
                                              │ (N:1)                │ (N:1)
                                              ↓                      ↓
                                      Utilisateur (formateur)   Utilisateur (apprenant)

Document ──(N:1)──> SessionFormation
Document ──(N:1)──> Utilisateur
\end{verbatim}

\subsection{Explication des relations}

\textbf{Entreprise ↔ DemandeFormation}
\begin{itemize}
    \item Une entreprise peut faire plusieurs demandes de formation
    \item Chaque demande est liée à une entreprise cliente
\end{itemize}

\textbf{Formation ↔ SessionFormation}
\begin{itemize}
    \item Une formation peut avoir plusieurs sessions (différentes dates/lieux)
    \item Chaque session est liée à une formation spécifique
\end{itemize}

\textbf{SessionFormation ↔ Utilisateur (formateur)}
\begin{itemize}
    \item Une session est animée par un formateur (utilisateur de type FORMATEUR)
    \item Un formateur peut animer plusieurs sessions
\end{itemize}

\textbf{SessionFormation ↔ Inscription}
\begin{itemize}
    \item Une session peut avoir plusieurs inscriptions
    \item Chaque inscription est liée à une session spécifique
\end{itemize}

\textbf{Inscription ↔ Utilisateur (apprenant)}
\begin{itemize}
    \item Un apprenant (utilisateur de type APPRENANT) peut s'inscrire à plusieurs sessions
    \item Chaque inscription est liée à un apprenant
\end{itemize}

\textbf{DemandeFormation ↔ Formation}
\begin{itemize}
    \item Une demande concerne une formation existante du catalogue
    \item Une formation peut faire l'objet de plusieurs demandes
\end{itemize}

\textbf{Document ↔ SessionFormation / Utilisateur}
\begin{itemize}
    \item Un document peut être associé à une session (ex: attestation de présence)
    \item Un document peut être associé à un utilisateur (ex: contrat)
\end{itemize}

\subsection{Flux métier principal}

\begin{enumerate}
    \item \textbf{Demande} : Une entreprise fait une demande pour une formation
    \item \textbf{Planification} : Une session est créée pour cette formation avec un formateur
    \item \textbf{Inscription} : Les apprenants s'inscrivent à la session
    \item \textbf{Réalisation} : La session se déroule (statut passe de PLANIFIEE à EN\_COURS puis TERMINEE)
    \item \textbf{Documents} : Les documents sont générés (attestations, factures, etc.)
\end{enumerate}

\section{Sécurité et authentification}

L'application utilise Keycloak pour la gestion de l'authentification et des autorisations :
\begin{itemize}
    \item Connexion sécurisée des utilisateurs
    \item Gestion des rôles et des permissions
    \item Protection des données sensibles
\end{itemize}

\begin{center}
\framebox{\parbox{0.8\textwidth}{\centering \textbf{[EMPLACEMENT SCREENSHOT - Interface Keycloak]} \\ \textit{Insérer ici une capture de l'interface Keycloak montrant la gestion des utilisateurs/rôles}}}
\end{center}

\section{Documentation de l'API}

Une interface Swagger a été mise en place pour documenter et tester l'API :
\begin{itemize}
    \item Documentation interactive de tous les endpoints
    \item Possibilité de tester les requêtes directement
    \item Accès via : \url{http://localhost:8080/swagger-ui.html}
\end{itemize}

\begin{center}
\framebox{\parbox{0.8\textwidth}{\centering \textbf{[EMPLACEMENT SCREENSHOT - Interface Swagger]} \\ \textit{Insérer ici une capture de l'interface Swagger UI avec les endpoints}}}
\end{center}

\section{Tests API avec Postman}

Les fonctionnalités ont été testées avec Postman :
\begin{itemize}
    \item Collection de requêtes pour chaque module
    \item Vérification des réponses de l'API
    \item Tests des différents scénarios
\end{itemize}

\begin{center}
\framebox{\parbox{0.8\textwidth}{\centering \textbf{[EMPLACEMENT SCREENSHOT - Postman - Collection de tests]} \\ \textit{Insérer ici une capture de Postman montrant la collection de requêtes}}}
\end{center}

\begin{center}
\framebox{\parbox{0.8\textwidth}{\centering \textbf{[EMPLACEMENT SCREENSHOT - Postman - Test d'un endpoint]} \\ \textit{Insérer ici une capture de Postman montrant un test réussi d'un endpoint API}}}
\end{center}

\section{Infrastructure et Docker}

L'application utilise Docker Compose pour l'orchestration des services :
\begin{itemize}
    \item Conteneur PostgreSQL pour la base de données
    \item Conteneur Keycloak pour l'authentification
    \item Facilité de déploiement et de reproduction de l'environnement
\end{itemize}

\begin{center}
\framebox{\parbox{0.8\textwidth}{\centering \textbf{[EMPLACEMENT SCREENSHOT - Docker Compose]} \\ \textit{Insérer ici une capture des conteneurs Docker en cours d'exécution}}}
\end{center}

\section{État d'avancement}

\subsection*{Fonctionnalités terminées}
\begin{itemize}
    \item Gestion complète des utilisateurs
    \item Gestion des entreprises
    \item Gestion des formations et sessions
    \item Gestion des inscriptions
    \item Gestion des demandes de formation
    \item Gestion des documents
    \item Dashboard de synthèse
    \item Authentification et sécurité
    \item Documentation API (Swagger)
\end{itemize}

\subsection*{En cours}
\begin{itemize}
    \item Tests finaux avec Postman
\end{itemize}

\subsection*{À faire}
\begin{itemize}
    \item Frontend, Ai, etc...
\end{itemize}

\section{Prochaines étapes}

\begin{itemize}
    \item Finalisation des tests
    \item Entamer la partie frontend avec Angular
\end{itemize}

\section*{Annexes}

\subsection*{Accès à l'application}

\begin{itemize}
    \item \textbf{Interface Swagger} : \url{http://localhost:8080/swagger-ui.html}
    \item \textbf{Port de l'application} : 8080
\end{itemize}

\end{document}
